
\subsubsection{A.1 Hypothesis Testing Summary}

\begin{table*}[htbp]
\centering
\resizebox{\dimexpr 2\columnwidth+\columnsep\relax}{!}{%
\begin{tabular}{lllll}
\toprule
Hypothesis & Type & Gate & Result & Key Metrics \\
\midrule
H-E1 & Efficacy & MUST\_WORK & PASS & 62.9\% discharge, 5.7 iterations, 100\% improvement \\
H-E2 & Efficacy & MUST\_WORK & PASS & 100\% coverage across 33 error categories \\
H-M1 & Mechanism & MUST\_WORK & PASS & $\beta =12.49, R^{2}=0.89$, p<$10^{-50}$ \\
H-M2 & Mechanism & SHOULD\_WORK & FAIL & Staged 57.2\% vs Complete 60.3\%, p=0.158 \\
H-M3 & Mechanism & MUST\_WORK & PASS & 15.1\% mean degradation, 6/6 pairs pass \\
H-C1 & Control & MUST\_WORK & PASS & 71.4\% vs 60.8\%, 10.7pp gap, p<0.0001 \\
H-C2 & Control & MUST\_WORK & PASS & 63.3\% vs 42\% threshold, 105.6\% of gold \\
\bottomrule
\end{tabular}
}
\end{table*}

\subsubsection{A.2 Semantic Primitive Definitions}

Each primitive includes semantic meaning, proof obligation type, matching keywords, and cross-verifier examples. The 8-primitive taxonomy achieves 100\% coverage:

\begin{itemize}
\item \textbf{MISSING\_PRECONDITION} (3 categories, 9.1\%): Entry condition under-specification
\item \textbf{POSTCONDITION\_FAILURE} (8 categories, 24.2\%): Exit guarantee under-specification
\item \textbf{LOOP\_INVARIANT\_VIOLATION} (6 categories, 18.2\%): Inductive invariant under-specification
\item \textbf{BOUNDS\_CHECK\_FAILURE} (3 categories, 9.1\%): Array/memory safety
\item \textbf{ARITHMETIC\_OVERFLOW} (5 categories, 15.2\%): Numeric safety including division-by-zero
\item \textbf{NULL\_DEREFERENCE} (3 categories, 9.1\%): Pointer safety
\item \textbf{TERMINATION\_FAILURE} (3 categories, 9.1\%): Liveness violations
\item \textbf{TYPE\_MISMATCH} (2 categories, 6.1\%): Type system violations
\end{itemize}

All primitives are semantically necessary, with no redundant categories.

\subsubsection{A.3 Mock Validation Details}

Mock validation used stochastic discharge rates to simulate verifier behavior for proof-of-concept experiments. This approach enables mechanism validation while controlling for real-world SMT solver variability and ensuring experimental reproducibility.

\textbf{RawError vs FullStructured Feedback Examples:}

\begin{table}[htbp]
\centering
\resizebox{\columnwidth}{!}{%
\begin{tabular}{lll}
\toprule
Program & RawError Feedback & FullStructured Feedback \\
\midrule
binary\_search & ''Postcondition may not hold at line 42'' & **Witness:** x=5, arr=[1,3,7,9], index=-1 violates ensures clause<br>**Structure:** POSTCONDITION\_FAILURE<br>**Dependency:** Loop invariant too weak $\rightarrow$ postcondition unprovable \\
find\_max & ''Assertion might fail at line 28'' & **Witness:** arr=[3, 5, 2], max=3 but arr[1]=5<br>**Structure:** LOOP\_INVARIANT\_VIOLATION<br>**Dependency:** Invariant doesn't track maximum across full traversed range \\
array\_copy & ''Precondition violation at line 15'' & **Witness:** src=NULL, len=10<br>**Structure:** MISSING\_PRECONDITION<br>**Dependency:** Requires valid\_read(src, len) precondition \\
\bottomrule
\end{tabular}
}
\end{table}

The 31.9\% RawError discharge rate provides a baseline representing unstructured-feedback performance.
